% !TeX root = ../CV note.tex
\section{Figure Detection}

\subsection{Edge Detection}

\paragraph{Mathematical Principle and Matrix Discretization}

边缘是图像中亮度发生突变（不连续）之处：表面法向、深度、颜色、光照的不连续都会形成边缘，而边缘携带了图像中大部分语义与形状信息，因此它是最先被提取的特征。数学上把灰度图视作二维函数 $f(x,y)$，"突变"就是函数在某处变化剧烈，用导数来度量。把沿 $x$、$y$ 的偏导数堆叠起来即得梯度
\[
\nabla f=\begin{bmatrix} f_x \\ f_y \end{bmatrix}
=\begin{bmatrix} \partial f/\partial x \\ \partial f/\partial y \end{bmatrix},
\qquad
|\nabla f|=\sqrt{f_x^2+f_y^2},
\qquad
\theta=\arctan\frac{f_y}{f_x}.
\]
在一维剖面上，边缘正对应一阶导数的极值；到二维就是梯度幅值 $|\nabla f|$ 的大值，$\theta$ 给出边缘的法向。

图像是离散采样的，导数只能由差商近似。从泰勒展开 $f(x+h)=f(x)+f'(x)h+\frac{f''(x)}{2!}h^2+\cdots$ 出发，取 $h=1$ 得前向差分 $f_x\approx f(x+1,y)-f(x,y)$；若把 $x+1$ 与 $x-1$ 两处的展开相减，二阶项相消，得到中心差分
\[
f_x\approx\frac{f(x+1,y)-f(x-1,y)}{2},
\]
它一步就有 $O(h^2)$ 精度（前向差分只有 $O(h)$），因此实际中一律采用中心差分。

真正让导数能在计算机上"跑起来"的是矩阵离散化：差商不过是相邻像素的加权和，因而可以写成一个卷积核。一维情形是 $[-1\ \ 1]$（前向）与 $[-1\ \ 0\ \ 1]$（中心）；二维时把一维核沿另一方向复制，再叠加平滑权，就得到
\[
H_x^{\text{Prewitt}}=\begin{bmatrix}-1&0&1\\-1&0&1\\-1&0&1\end{bmatrix},\quad
H_y^{\text{Prewitt}}=\begin{bmatrix}1&1&1\\0&0&0\\-1&-1&-1\end{bmatrix},
\]
\[
H_x^{\text{Sobel}}=\begin{bmatrix}-1&0&1\\-2&0&2\\-1&0&1\end{bmatrix},\quad
H_y^{\text{Sobel}}=\begin{bmatrix}1&2&1\\0&0&0\\-1&-2&-1\end{bmatrix}.
\]
于是求导退化为 3.2 节讲过的卷积滤波：$f_x=f*H_x$，$f_y=f*H_y$。核中的 $\pm1$ 在做差商、$0$ 在取中心、权重 $2$ 与 $1$ 在做平滑——Prewitt 与 Sobel 的唯一差别就在平滑权：Sobel 把中间行加倍，相当于在求导方向上再叠一个 $[1,2,1]$ 的二项式核，对噪声更稳健。

\begin{figure}[H]
    \centering
    \includegraphics[width=0.85\textwidth]{/home/simson/latex/CSAI note/CV/figs/edge-prewitt-sobel-kernels.png}
    \caption{Prewitt 与 Sobel 的导数核：水平核（对 $x$ 求导）与垂直核（对 $y$ 求导）都是 3$\times$3 矩阵，即离散化后的偏导算子}
\end{figure}

把 $|\nabla f|$ 阈值化即可得到边缘图（Sobel 边缘检测器），但阈值很脆弱、很难选：偏高会漏掉弱边缘，偏低又会把噪声当成边缘——这正是后面 Canny 用非极大值抑制与双阈值要解决的问题。

\paragraph{Noise Suppression and the Gaussian Derivative}

噪声对导数是致命的。设观测图像为 $I_{i,j}=T_{i,j}+\epsilon_{i,j}$，其中 $T$ 是真实图像，$\epsilon_{i,j}\sim N(0,\sigma^2)$ 是各像素独立同分布的噪声。用中心差分求导时，噪声被一并相减：
\[
D_{i,j}=(T_{i,j+1}+\epsilon_{i,j+1})-(T_{i,j-1}+\epsilon_{i,j-1})
=(T_{i,j+1}-T_{i,j-1})+(\epsilon_{i,j+1}-\epsilon_{i,j-1}).
\]
后一项是两个独立高斯之差，仍是高斯，但方差相加：$\epsilon_{i,j}-\epsilon_{k,l}\sim N(0,2\sigma^2)$，即\textbf{方差翻倍}，标准差被放大 $\sqrt2$ 倍。平坦区域里真实差分近似为零，导数图便由纯噪声主导，真正的边缘被完全淹没。

解决办法是"先平滑、再求导"：先用高斯核 $g$ 平滑图像再做微分，即求 $\frac{d}{dx}(f*g)$。由卷积与微分可交换，
\[
\frac{d}{dx}(f*g)=f*\frac{dg}{dx},
\]
所以两趟卷积能合并成一趟：只需先把高斯核求导得到"高斯导数"核 $\frac{dg}{dx}$（一维上是一侧为正、一侧为负的反对称波形），再与图像卷积一次，就同时完成了去噪与求导。二维下对应两个核 $G_x=\partial G_\sigma/\partial x$ 与 $G_y=\partial G_\sigma/\partial y$，按网格离散化就是对采样点直接求偏导：
\[
G_x[i,j]=-\frac{i}{\sigma^2}e^{-\frac{i^2+j^2}{2\sigma^2}},\qquad
G_y[i,j]=-\frac{j}{\sigma^2}e^{-\frac{i^2+j^2}{2\sigma^2}},
\]
于是 $f_x=f*G_x$、$f_y=f*G_y$，一个水平、一个垂直。

\begin{figure}[H]
    \centering
    \includegraphics[width=0.95\textwidth]{/home/simson/latex/CSAI note/CV/figs/gaussian-derivative-1d.png}
    \caption{高斯导数：$\frac{d}{dx}(f*g)=f*\frac{dg}{dx}$ 把去噪与求导合并为一趟卷积；右侧是二维的 $\frac{\partial G_\sigma}{\partial x}$ 与 $\frac{\partial G_\sigma}{\partial y}$}
\end{figure}

代价是尺度上的两难：$\sigma$ 越大去噪越彻底，但边缘同时被抹平，导数峰展宽、定位精度下降（课件分别用 1、3、7 像素的高斯核：越小噪声残留越多，越大边缘越糊）。"去噪"与"保边定位"无法兼得，$\sigma$ 只能折中——这也解释了为什么 Canny 的第一步（灰度化 + 高斯滤波）与 Harris 角点的预处理都要先做高斯平滑。

\paragraph{Canny Edge Detector}

Canny 检测器把前面几节的工具串成一条六步流水线，目标是得到"细、连续、定位准"的边缘：

\begin{enumerate}[label=\arabic*.]
    \item \textbf{预处理：灰度化 + 高斯模糊}。先把彩色图转成灰度图（边缘本是由亮度突变定义的），再用高斯核平滑以抑制噪声——这正是上一段的动机：求导会放大噪声，所以必须先平滑。
    \item \textbf{中心差分求梯度图}。用 Sobel 的中心差分算出 $G_x$、$G_y$（而不用前向差分，因为中心差分精度更高）。$G_x$ 突出竖直边缘、$G_y$ 突出水平边缘。
    \item \textbf{计算梯度幅值}。$|\nabla f|=\sqrt{G_x^2+G_y^2}$，即边缘强度图。
    \item \textbf{非极大值抑制（Non-Maximum Suppression）}。幅值图上的边缘是粗的：边缘附近一整条带上的像素幅值都很大。做法是遍历幅值矩阵，只保留沿\textbf{梯度方向}的局部极大值。对像素 $q$，沿梯度方向取它两侧的邻居 $p$ 与 $r$；梯度方向一般不落在整数格点上，故 $r$ 由相邻两点线性插值得到 $r=\alpha b+(1-\alpha)a$。当且仅当 $q>p$ 且 $q>r$ 时保留 $q$，否则置零——一条粗边就被削成单像素宽。
    \item \textbf{双阈值}。设高阈值（如 $0.7$）与低阈值（如 $0.3$）：高于高阈值的像素确定是边缘（强边缘），低于低阈值的一定不是，介于两者之间的是\textbf{弱边缘}，留给下一步判定。用双阈值取代单阈值，正是为了化解上一段末尾说的"阈值难选"：单阈值偏高会漏掉弱边缘，偏低又会引入噪声。
    \item \textbf{滞后边缘跟踪（Edge Tracking by Hysteresis）}。与强边缘相连的弱边缘判为真实边缘而保留，不与任何强边缘相连的弱边缘则丢弃。这一步让边缘在弱对比处仍能连续延伸，同时排除孤立的噪声响应。
\end{enumerate}

\begin{figure}[H]
    \centering
    \includegraphics[width=0.95\textwidth]{/home/simson/latex/CSAI note/CV/figs/canny-non-maximum-suppression.png}
    \caption{非极大值抑制：沿梯度方向比较 $q$ 与两侧的 $p$、$r$（$r$ 由 $a$、$b$ 线性插值），只保留局部极大值，把粗边缘削成单像素宽}
\end{figure}

\paragraph{Applications of the Canny Detector}

Canny 给出的单像素、连续的边缘图本身就是物体的"线稿"，因此常被当作下游任务的条件输入：

\begin{itemize}
    \item \textbf{草图生成与三维形状生成}。Locally Attentional SDF Diffusion for Controllable 3D Shape Generation（SIGGRAPH 2023）用 Canny 从渲染图中提取线稿，再以线稿为条件生成可控的三维形状。
    \item \textbf{文本到图像的附加控制}。Adding Conditional Control to Text-to-Image Diffusion Models 中用 Canny 边缘图作为空间条件注入扩散模型，于是"一张线稿 + 一句文本描述"就能精确控制生成图像的构图。
    \item \textbf{通用轮廓提取}。图像编辑的选区引导、工业检测的缺陷轮廓、多视图几何中的特征线提取等，都把 Canny 当作默认的边缘算子。
\end{itemize}

\begin{figure}[H]
    \centering
    \includegraphics[width=0.95\textwidth]{/home/simson/latex/CSAI note/CV/figs/canny-application-sketch3d.png}
    \caption{应用一：用 Canny 从渲染图提取线稿，作为三维形状生成的可控条件（Locally Attentional SDF Diffusion, SIGGRAPH 2023）}
\end{figure}

\begin{figure}[H]
    \centering
    \includegraphics[width=0.95\textwidth]{/home/simson/latex/CSAI note/CV/figs/canny-application-controlnet.png}
    \caption{应用二：Canny 线稿作为扩散模型的空间控制条件，同一张线稿配不同文本可生成不同风格的图像（ControlNet / Adding Conditional Control to Text-to-Image Diffusion Models）}
\end{figure}

不过课件最后提醒：\textbf{边缘检测至今仍是一个未解决的问题}。人眼能把图中的物体轮廓干净地分割出来，而梯度幅值图在物体内部与背景处的杂散响应同样强烈，二者相差甚远。这也正是接下来要从边缘转向更能刻画局部结构的特征——角点与 blob——的原因。

